\documentclass[11pt]{article}
\usepackage[a4paper,margin=25mm]{geometry}
\usepackage{amsmath,amssymb,amsthm,mathtools,booktabs,hyperref,xurl,microtype}
\setlength{\emergencystretch}{3em}
\hypersetup{hidelinks}
\newtheorem{theorem}{Theorem}
\newtheorem{lemma}[theorem]{Lemma}
\newcommand{\Q}{\mathbb Q}\newcommand{\Z}{\mathbb Z}
\newcommand{\re}{\operatorname{re}}
\title{Directly scaled private points:\\a 645-variable ordinary-abc quartic transfer}
\author{\(\Z\Q\) Diophantine Project}
\date{24 November 2026}
\begin{document}\maketitle

\begin{abstract}
One additional common product lets the private ordinate ratio and chord
remain input-scaled. This removes the unscaled ordinate-ratio field and
unscaled half-trace used in the preceding common-point presentation.
The resulting complete table has 645 rational unknowns and 564 quadratic
residuals. We give a polynomial lift directly to the frozen 674/593
odd-Pell predecessor, valid on every rational zero including zero-input
fibers. Completeness uses the same common-numerator arithmetic theorem.
All retained source premises, ordinary rational abc, and the ineffective
defining coefficient remain required. This is a fresh construction,
not a recovery of the unavailable earlier 649-coordinate artifact.
\end{abstract}

\section{The result and its arithmetic scope}

The common-numerator theorem~\cite{common} supplies a common physical
point \(U\), a common reference point, and private differences for every
tuple of nonzero integer inputs. It retains the elliptic group,
comparison, trace-index, finite prime-pattern, local preparation, real
recurrence and pole-realization premises of its ordinary-abc source.
Ordinary rational abc is also assumed, and its sufficiently large odd
defining coefficient remains ineffective. We make no new arithmetic
assertion about those premises.

\begin{theorem}\label{thm:main}
Under exactly those retained premises, the simultaneous ordinary-abc
integer interface has a quadratic presentation with
\[
        43k+24\quad\text{rational coordinates},\qquad
        38k+19\quad\text{quadratic residuals},
\]
including its \(k\) distinguished inputs. Every rational zero has
integral inputs. Every nonzero integer tuple has a rational witness.
For the thirteen certified inputs in the odd-Pell source~\cite{odd},
the complete table has 645 rational unknowns and 564 quadratic residuals;
its instance projection agrees with that of the 674-variable predecessor.
The sum of squared residuals has degree exactly four.
\end{theorem}

The source instance is a coefficient. All thirteen integer predicates,
the externally cited Jones coding theorem, and the source's canonical
positive completion are retained. There is no universality premise for
the historically emended source and no unconditional undecidability
claim over \(\Q\). The numerical comparison with the user's earlier
649-coordinate checkpoint is a comparison of counts only: its original
artifact was not available, so no assertion about its contents is made.

\section{The quadratic chart}

Let \(K=\Q(\rho)\), \(\rho^2=2\), and \(\Delta=11-6\rho\).
Every field coordinate has two rational components. For
\(a+b\rho\), put \(\re(a+b\rho)=a\); a bar denotes conjugation.
The common point \(U\) on \(y^2=x^3-\Delta x\) has fields
\(S=y(U)/x(U)\) and \(F=1/x(U)\). Retain shared fields
\(S_2,S_R,S_F,\sigma\), with
\begin{equation}\label{eq:shared}
 S_2=S^2,\quad S_R=SR,\quad S_F=SF,\quad
 FS_2+\Delta F^2=1,\quad uF=vJ,\quad w\sigma=v.
\end{equation}
Here \(J\) is the previously shared reciprocal product; its retained
square field \(K_J=J^2\) and its two defining rows are unchanged.
The shared reference fields are \(R=S_H\), \(T_0=S_{H-R_0}\), \(v=v_H\),
where \(R_0=(1+\rho,2\rho)\), and
\[
 u=(\rho-1)(v-(4-2\rho)-R),\qquad
 w=(\rho-1)(v-(4-2\rho)+T_0).
\]
The existing reference chart makes \(u,w,v\) nonzero. Thus
\(\sigma=v/w\) exists in the completeness construction. All six
field equations in~\eqref{eq:shared} are quadratic.

For one distinguished input \(X\), keep the predecessor's registers
\(X_2=X^2\), \(X_4=X_2^2\), chart units \(c\epsilon=1\),
half-trace \(a\), input-scaled ordinate field \(T\), and input-scaled
chord field \(H\). Add two private fields \(G,Z\) and one private
scalar \(\omega\). The ten private rational equations are
\begin{align}
 (H+T)F+XS_F-GH&=0, \label{eq:Aquad}\\
 \Delta X_2(F+G)-SZ+TH+T^2&=0, \label{eq:Bquad}\\
 Z-XH&=0, & \omega-cX_2&=0, \label{eq:Zweight}\\
 a-\re(H^2+SZ+HT)&=0, \label{eq:trace}\\
 \sigma\overline L-\omega G&=0. \label{eq:L}
\end{align}
The field \(L\) is the private reciprocal-difference field already
retained in~\cite{common}. No new unscaled ordinate or half-trace
coordinate is introduced.

For the retained comparison fields \(P,Q\), replace the first
comparison equation by
\begin{equation}\label{eq:compare}
       2(X_2S_2-2XS_R+R^2)-P-cQ=0.
\end{equation}
Its second comparison equation uses the retained \(T\) and is
unchanged. Every other contact, detector, pole, height, source and
power-coding equation remains unchanged after coordinate reindexing.

\section{A polynomial lift with no division by the input}

Write
\begin{align*}
 A&=(H+XS+T)F-GH,\\
 B&=\Delta X^2(F+G)-(XS-T)H+T^2,\\
 C&=FS^2+\Delta F^2-1.
\end{align*}
The quadratic equations above imply \(A=B=C=0\). The following
identities are polynomial identities in unrestricted field variables:
\begin{align}
 GH(XS-T)-\Delta X^2FG-X^2
   &=X^2C-FB-(XS-T)A, \label{eq:scaleid}\\
 GT^2+\Delta X^2G^2-X^2
   &=X^2C+(G-F)B-(XS-T)A. \label{eq:curveid}
\end{align}
If \(X\ne0\), the second identity and the shared curve recover the
private curve and its nonzero reciprocal abscissa. The soundness proof
below does not divide by \(X\), and uses neither of those conclusions
on the \(X=0\) fiber.

\begin{lemma}\label{lem:lift}
Every rational zero of the new complete table lifts to a zero of the
complete 674/593 predecessor with the same distinguished inputs and
source instance, by polynomial reconstruction of the deleted coordinates.
\end{lemma}
\begin{proof}
Use, block by block,
\begin{align*}
 s^{\rm old}&=XS,& F^{\rm old}&=X_2F,&
 G^{\rm old}&=\omega G,\\
 t^{\rm old}&=T,& h^{\rm old}&=H,&
 d^{\rm old}&=\omega GH,
\end{align*}
and set the deleted common field \(E^{\rm old}=vJ\).
The already-retained equations \(X_2=X^2\), \(X_4=X^4\),
\(c\epsilon=1\), and \(\omega=cX_2\) identify the old pair
first-product, second-product and secant residuals with \(X^2A\),
zero and \(-B\). The old scale residual becomes
\[
 cX^2\bigl(GH(XS-T)-\Delta X^2FG-X^2\bigr),
\]
which vanishes by~\eqref{eq:scaleid}. In particular this argument
continues to hold identically at \(X=0\).

The old difference half-trace is recovered because
\[
 H(H+XS+T)-(H^2+SZ+HT)=S(XH-Z).
\]
Its real component vanishes by~\eqref{eq:Zweight}, and the retained
\(a\) satisfies~\eqref{eq:trace}. The first reciprocal-product
residual is \(X^2(uF-vJ)=0\); the second follows from
\[
 v\overline L-w\omega G
       =w(\sigma\overline L-\omega G)-(w\sigma-v)\overline L.
\]
The common \(E\)-definition is automatic. Finally,
\(X_2S_2-2XS_R+R^2=(XS-R)^2\), so~\eqref{eq:compare}
is exactly the old first comparison relation after substitution.
Every other old row is unchanged and contains no deleted coordinate.
No division by an input or by a new private field has occurred.
\end{proof}

This is a direct lift to~\cite{odd}. At zero inputs the retained
\(H,T\) need not be input multiples, so it is not claimed that every
new auxiliary tuple lifts to the intermediate 656-coordinate
presentation~\cite{unscaled}. The common source instance projection
is the relevant equivalence.

\section{Completeness and exact accounting}

For nonzero integer inputs, the common-numerator theorem~\cite{common}
chooses one physical \(U\), one reference point, and private points
\(V_i=U-[X_i]R_0\). It preserves the fixed differences, comparison
prime conditions to arbitrary local accuracy, fixed-place and real
conditions, contact locations and pole witnesses. We impose no new
condition on that choice. For each input set
\[
 T_i=X_iS_{V_i},\qquad H_i=X_ih_i,\qquad G_i=f(V_i),
 \qquad Z_i=X_iH_i,\quad \omega_i=c_iX_i^2,
\]
where \(h_i\) is the unscaled translated-secant value. Its original
formulas include tangents and zero ordinates~\cite{unscaled,secant}.
The unscaled secant identities, multiplied by \(X_i\) and \(X_i^2\),
give~\eqref{eq:Aquad}--\eqref{eq:Bquad}. The half-trace is precisely
the old scaled output. Take \(\sigma=v/w\), \(S_F=SF\),
\(S_2=S^2\), and \(S_R=SR\); the old reciprocal relation divided
by the reference unit \(w\) gives~\eqref{eq:L}. This proves
completeness without altering the pole data or arithmetic premises.

The previous common-reciprocal tuple presentation has \(46k+14\)
coordinates and \(41k+9\) rows. Remove eight old rational point
coordinates per input, and add five private rational coordinates.
Remove the two coordinates of \(E\), and add twelve shared rational
coordinates. Thus
\[
             (46k+14)-8k+5k-2+12=43k+24.
\]
Thirteen old rows per input are replaced by ten. Two comparison rows
per input are changed but not added. Replacing the two old
\(E\)-definitions by twelve common rows gives
\[
                       (41k+9)-3k-2+12=38k+19.
\]
At \(k=13\) this saves 29 coordinates and 29 rows from the odd-Pell
674/593 table, giving 645/564. It also saves eleven coordinates and
rows from the 656/575 presentation of~\cite{unscaled}. Every residual
is quadratic over the retained integer parameter ring. The unchanged
input-square residual has nonzero quadratic part, and squares of
rational leading quadratic forms cannot cancel, so the sum of squared
residuals has degree exactly four. Lemma~\ref{lem:lift} and the
completeness construction prove Theorem~\ref{thm:main}.

\section{Reproducible verification}

The directory contains the complete table, deterministic compiler,
independent verifier, metadata certificate and independent audit receipt.
Run from the repository root:
\begin{center}\small
\texttt{python -B}\enspace\nolinkurl{undecidability/abc_scaled_point_v106/build_scaled_point.py}\\
\texttt{python -B}\enspace\nolinkurl{undecidability/abc_scaled_point_v106/verify_scaled_point.py}.
\end{center}
The independent verifier does not import the compiler. It checks 197
old row interpretations, all 168 new or changed rows, all 396 unchanged
rows, every occurrence of all 106 removed coordinates, and nine
universal field identities. It also checks the coefficient domain,
full support, degree, counts and table hashes. These checks certify the
finite algebraic construction; they do not certify the retained
arithmetic premises or provide an effective sufficient coefficient.

\small
\begin{thebibliography}{9}\raggedright
\setlength{\itemsep}{2pt}
\bibitem{common} \(\Z\Q\) Diophantine Project,
\emph{A common numerator compresses reciprocal products in the ordinary-abc
quartic transfer}, 24 November 2026,
\nolinkurl{undecidability/abc_common_numerator_v103/common_numerator.tex}.
\bibitem{odd} Project, \emph{An odd second-Pell quotient removes one inverse},
24 November 2026,
\nolinkurl{undecidability/odd_pell_quotient_v104/odd_pell_quotient.tex};
complete predecessor
\nolinkurl{undecidability/odd_pell_quotient_v104/odd_quotient_abc_quadrics.json}.
\bibitem{unscaled} Project, \emph{Sharing the physical numerator chart in
the ordinary-abc quartic transfer}, 24 November 2026,
\nolinkurl{undecidability/abc_point_rechart_v105/common_point_rechart.tex}.
\bibitem{secant} Project, \emph{A quartic in 68 variables from translated
secants}, 17 November 2026,
\nolinkurl{papers/quartic_existence/quartic_abc68_20261117_v65.tex}.
\end{thebibliography}
\end{document}
